\chapter{Thiết kế kiểm soát bảo mật}

\section{Xác thực và phiên}

Endpoint \texttt{/auth/login} xác thực username và password qua repository.
Khi thành công, backend tạo session token và trả về thông tin user cùng role;
khi thất bại, backend ghi sự kiện audit nhưng không tiết lộ tài khoản nào tồn
tại. Frontend đặt token trong \texttt{sessionStorage}, vì vậy mỗi tab có phiên
riêng phục vụ kịch bản demo đồng thời Admin, Teller và Controller. Các endpoint
có bảo vệ dùng Bearer token; token thiếu, sai hoặc hết hiệu lực nhận HTTP 401.

\section{Phân quyền tại backend}

FastAPI đóng gói kiểm tra quyền qua dependency \texttt{require\_permission}.
Dependency này gọi \texttt{RBACRepository.has\_permission}, ghi một audit event
cho cả quyết định allowed và denied, sau đó chỉ cho handler chạy khi quyền hợp
lệ. Ví dụ, quản lý user đòi \texttt{users:manage}, xem audit log đòi
\texttt{audit\_logs:read}, còn tạo giao dịch đòi \texttt{transactions:create}.
Do enforcement nằm ở API, người dùng không thể vượt quyền bằng cách mở URL hay
gửi request trực tiếp ngoài giao diện.

\section{Separation of duties}

Chính sách phê duyệt có hai chế độ. \emph{Baseline} giữ một bypass có chủ đích
để Teller có thể duyệt, tạo đối chứng cho bài trình diễn. \emph{RBAC} loại bỏ
bypass này: request xem hoặc phê duyệt hồ sơ gọi policy enforcement, chỉ cho
phép role có \texttt{transactions:approve}. Teller truy cập URL
\texttt{/approvals/\{reference\}} sau khi chuyển sang RBAC nhận HTTP 403; giao
dịch vẫn pending và không phát sinh approver.

\section{Audit trail và tính nhất quán}

Mỗi thao tác có ý nghĩa bảo mật được lưu kèm actor, resource, action, outcome,
thời điểm và, khi có, mã giao dịch. Với phê duyệt, repository xử lý trạng thái
giao dịch theo thao tác nguyên tử. Nếu giao dịch đã được xử lý, API trả HTTP
409 và audit outcome \texttt{conflict}; do đó kết quả không phụ thuộc vào việc
người dùng bấm lại nút hoặc nhiều request đến sát nhau. Audit log không thay
thế request/gateway log chứa chỉ báo SQL Injection, nhưng adapter
\texttt{export-audit-events} đã cho phép bản ghi xác thực và phân quyền đi qua
pipeline chung. Thực nghiệm đầu-cuối tạo chuỗi đăng nhập thất bại và một thao
tác bị từ chối, xuất audit rồi quan sát hai luật tương ứng tạo alert.

Tính toàn vẹn của audit trail trong PoC phụ thuộc vào database và quyền file
cục bộ. Hệ thống chưa có write-once storage, hash chain, đồng bộ thời gian tin
cậy hay chuyển log sang collector độc lập. Đây là ranh giới quan trọng: audit
trail hiện tại phù hợp minh họa trách nhiệm giải trình của ứng dụng, chưa phải
bằng chứng pháp chứng chống sửa đổi.
